\centerline{\bf \Huge Полная комбинаторика}
\title{Полная комбинаторика}


1.1 Сколькими способами можно разбить 15 человек на три команды по пять человек в каждой?

1.2 Сколькими способами можно выбрать из 15 человек две команды по пять человек в каждой?

2. Переплетчик должен переплести 12 одинаковых книг в красный, зелёный или синий переплеты. Сколькими способами он может это сделать?

3. Сколькими способами три человека могут разделить между собой шесть одинаковых яблок, один апельсин, одну сливу и один мандарин?

4. Сколькими способами четыре чёрных шара, четыре белых шара и четыре синих шара можно разложить в шесть различных ящиков?

5. Мария Ивановна покупает 16 шариков для Последнего звонка. В магазине есть шарики трёх цветов: синего, красного и зелёного. Сколько существует вариантов различных покупок 16 шариков, если Мария Ивановна хочет, чтобы шарики каждого цвета составляли не менее четверти от количества всех шариков?

6. Сколькими способами можно разделить на команды по 6 человек для игры в волейбол группу:
а) из 12; б) из 24 спортсменов?

7. Сколькими способами можно составить букет из 17 цветков, если в продаже имеются гвоздики, розы, гладиолусы, ирисы, тюльпаны и васильки?

8. Сколькими способами можно разбить 10 человек на две баскетбольные команды по 5 человек в каждой?

9. Имеется резинка и стеклянные шарики-бусины: четыре одинаковых красных, две одинаковых синих и две одинаковых зелёных. Нужно все восемь бусин нанизать на резинку последовательно, чтобы получился браслет. Сколько различных браслетов можно составить так, чтобы бусины одного цвета не оказались рядом? (Считайте, что застёжки нет, а узелок на резинке незаметен.)

10. В шахматном кружке занимаются 2 девочки и 7 мальчиков. Для участия в соревновании необходимо составить команду из четырёх человек, в которую обязательно должна входить хотя бы одна девочка. Сколькими способами это можно сделать?

11. Из 12 девушек и 10 юношей выбирают команду, состоящую из пяти человек.
Сколькими способами можно выбрать эту команду так, чтобы в нее вошло не более трёх юношей?

\newpage 

12. Известно, что в выпуклом n-угольнике (n > 3) никакие три диагонали не проходят через одну точку.
Найдите число точек (отличных от вершины) пересечения пар диагоналей.

13. 30 человек голосуют по пяти предложениям. Сколькими способами могут распределиться голоса, если каждый голосует только за одно предложение и учитывается лишь количество голосов, поданных за каждое предложение?

14. Из двух математиков и десяти экономистов надо составить комиссию из восьми человек.

Сколькими способами можно составить комиссию, если в неё должен входить хотя бы один математик?

15. Сколько существует шестизначных чисел, у которых по три чётных и нечётных цифры?